\documentclass[11pt, a4paper]{article}

% ── Packages ──────────────────────────────────────────────────────────────────
\usepackage[margin=1in]{geometry}
\usepackage{amsmath}
\usepackage{booktabs}
\usepackage{graphicx}
\usepackage{hyperref}
\usepackage{setspace}
\usepackage{parskip}
\usepackage{titlesec}
\usepackage{caption}
\usepackage{float}
\usepackage{lmodern}
\usepackage[T1]{fontenc}
\usepackage{microtype}
\usepackage{array}
\usepackage{xcolor}
\usepackage{tabularx}

% ── Formatting ─────────────────────────────────────────────────────────────────
\setstretch{1.15}
\titleformat{\section}{\normalfont\large\bfseries}{}{0em}{}[\titlerule]
\titleformat{\subsection}{\normalfont\normalsize\bfseries}{}{0em}{}
\titlespacing{\section}{0pt}{12pt}{5pt}
\titlespacing{\subsection}{0pt}{8pt}{3pt}
\captionsetup{font=small, labelfont=bf}
\hypersetup{colorlinks=true, linkcolor=black, citecolor=black, urlcolor=black}
\newcolumntype{C}[1]{>{\centering\arraybackslash}p{#1}}

\begin{document}

% ── Title Block ────────────────────────────────────────────────────────────────
\begin{center}
    {\ \textbf{Sovereign Debt Repricing under Geological Fragmentation in Central America}}\\[6pt]
    {\ \text{Note: This case utilized Gen AI tools for analysis but all ideas are original}}
\end{center}

\vspace{6pt}
\hrule
\vspace{2pt}
\hrule
\vspace{6pt}

% ── Bottom Line ────────────────────────────────────────────────────────────────
\noindent\colorbox{gray!12}{\begin{minipage}{\dimexpr\textwidth-2\fboxsep}
\vspace{4pt}
\textbf{Bottom Line.} If tomorrow, Central America drifted apart into islands, an interesting point of exploration would be the impact on sovereign borrowing costs. This study employs a cross-sectional regression across 22 emerging and frontier market sovereigns. We do not find a statistically significant globalized ``island premium''. However, this quantitatively reveals that markets likely price \textit{what kind of island} a country becomes, not island status generically. Central American nations would become the wrong kind. Applying the spread differential observed between comparables shows that the region's \$68 billion in external debt implies \textbf{\$200M--\$1.2B in additional annual interest costs}, central estimate $\sim$\textbf{\$680M per year}.
\vspace{4pt}
\end{minipage}}

% ══════════════════════════════════════════════════════════════════════════════
\section{Hypothesis}
% ══════════════════════════════════════════════════════════════════════════════

Island nations face very distinct risk profiles to their continental peers: greater hurricane and climate exposure, large food import dependency, higher trade costs overseas, and isolation from critical infrastructure. In this case, the fragmentation of Central America would sever the Pan-American Highway and destroy the SIEPAC regional power grid at the very least. Seven continental sovereigns would become island states overnight. Thus, it can be  theorized that this should manifest as a persistent premium in sovereign borrowing costs within the region.

% ══════════════════════════════════════════════════════════════════════════════
\section{Data and Methodology}
% ══════════════════════════════════════════════════════════════════════════════

The analysis utilizes a cross-sectional ordinary least-squares regression across 22 small-to-mid-sized emerging/frontier market sovereigns, averaged over 2013--2023 to explore structural pricing rather than cyclical pricing. The sample spans Central America (5 continental), the Caribbean (5 island), South America (4 continental), West and East Africa (5 continental), and the Indian Ocean and Pacific (3 island). Sovereigns were largely chosen based on availability of public data and comparability to Central America. The dependent variable being explored is the average sovereign bond spread over the US 10-year Treasury benchmark (2.32\% average, 2013--2023).

Three models estimated with increasing specifications for island conditions:

\begin{itemize}
    \item \textbf{Model 1 (Baseline):} Controls only: log GDP per capita, debt-to-GDP, CA \% of GDP
    \item \textbf{Model 2:} Controls + island binary dummy
    \item \textbf{Model 3 (Full):} Controls + island dummy + three natural disaster exposure score, logistics performance index as a remoteness proxy, and a food import dependency ratio
\end{itemize}

All models used HC3 heteroskedasticity-robust standard errors. Models 4 and 5 repeat Models 2 and 3 excluding Zambia and Ghana, which experienced sovereign debt restructurings in 2020 and 2023 respectively --- credit events unrelated to geography that distort the continental baseline.

% ══════════════════════════════════════════════════════════════════════════════
\section{Results}
% ══════════════════════════════════════════════════════════════════════════════

\subsection{Regression Output}

\begin{table}[H]
\centering
\caption{OLS Regression Results --- Island Dummy Coefficient}
\label{tab:regression}
\begin{tabular}{llC{1.1cm}C{1.1cm}C{2.2cm}C{1.6cm}}
\toprule
Model & Sample & N & R\textsuperscript{2} & Island Coef (bps) & $p$-value \\
\midrule
M2: Controls + dummy & Full                    & 22 & 0.469 & $-137$ & 0.400 \\
M3: Full model       & Full                    & 22 & 0.675 & $+110$ & 0.838 \\
M4: Controls + dummy & Clean & 20 & 0.383 & $-72$  & 0.519 \\
M5: Full model       & Clean                   & 20 & 0.628 & $+7$   & 0.983 \\
\end{tabular}
\end{table}

The island dummy is statistically insignificant across every specification, with $p$-values ranging from 0.40 to 0.98. Coefficient sign flips between M2/M4 and M3/M5 as the island-specific risk variables are added, reflecting high estimation uncertainty with 22 observations. Controls performed as expected. The full model R\textsuperscript{2} of 0.63--0.68 indicates the variable set explains meaningful variation overall, however the island dummy specifically cannot be identified.

\subsection{Why the Island Premium Is Not Identified}

Two structural features prevent clean identification. First, the islands with liquid sovereign bond markets --- Bahamas, Barbados, Mauritius, Seychelles, Trinidad and Tobago --- are wealthier than the continental comparison group. These are tourism and financial-hub economies whose spreads are structurally low not because of their geography but despite it. Second, the two continental outliers --- Zambia (2,168 bps) and Ghana (1,718 bps) --- inflate the continental average with default risk unrelated to geographic pricing. Excluding them brings the clean-sample continental average from 726 bps to 523 bps, narrowing the raw gap with islands at 403 bps. But even in the clean sample the dummy remains insignificant.

\medskip
\noindent This is not a failure, however. A general premium requires island sovereigns to trade wider than continental peers after controlling for fundamentals — but the islands that access bond markets are predominantly tourism and financial-hub economies whose geography is offset by hard-currency revenues and strong institutions. They trade tight \textit{despite} being islands. The premium only materialises for the kind of island Central America would become: mid-income, commodity-dependent, and newly deprived of overland trade infrastructure. Better insight into this premium would require a more intensive model and data, but nevertheless this study provides interesting insights into the relationship between continental and island sovereigns.

% ══════════════════════════════════════════════════════════════════════════════
\section{Finding: The Premium Depends on the Type of Island}
% ══════════════════════════════════════════════════════════════════════════════

Narrowing the comparison to within-region matched pairs --- Caribbean islands versus Central American continental nations at comparable income levels --- produces a relevant estimate.

\begin{table}[H]
\centering
\caption{Within-Region Matched Pairs: Sovereign Spreads (bps, avg.\ 2013--2023)}
\label{tab:matched}
\begin{tabular}{lC{2cm}lC{2.2cm}C{2cm}}
\toprule
Island & Spread (bps) & Comparable Continental & Spread (bps) & Differential \\
\midrule
Jamaica             & 488 & Guatemala   & 308 & $+180$ \\
Dominican Republic  & 548 & Honduras    & 528 & $+20$  \\
Barbados            & 588 & El Salvador & 558 & $+30$  \\
Trinidad \& Tobago  & 258 & Panama      & 248 & $+10$  \\
Bahamas             & 228 & Panama      & 248 & $-20$  \\
\midrule
\multicolumn{4}{l}{\text{Defensible range: 30--180 bps $|$ Central estimate: $\sim$100 bps}} & \\
\end{tabular}
\end{table}

% ══════════════════════════════════════════════════════════════════════════════
\section{Dollar Impact}
% ══════════════════════════════════════════════════════════════════════════════

Applied to the five available Central American sovereigns' aggregate external debt of approximately \$68 billion, a 100 bps structural spread widening implies roughly \$680 million in additional annual interest costs. At 180 bps the total is \$1.2 billion per year; at 30 bps it is \$200 million.

\begin{table}[H]
\centering
\caption{Estimated Annual Interest Cost Increase by Scenario}
\label{tab:dollar}
\begin{tabular}{lC{2.2cm}C{2.2cm}C{2.2cm}C{2.2cm}}
\toprule
Country & Ext.\ Debt (\$bn) & 30 bps (\$mn/yr) & 100 bps (\$mn/yr) & 180 bps (\$mn/yr) \\
\midrule
Guatemala     & 14.5 & 44  & 145 & 261   \\
Honduras      & 11.2 & 34  & 112 & 202   \\
El Salvador   &  8.8 & 26  &  88 & 158   \\
Costa Rica    & 13.4 & 40  & 134 & 241   \\
Panama        & 20.1 & 60  & 201 & 362   \\
\midrule
\text{Total} & \text{68.0} & \text{204} & \text{680} & \text{1,224} \\
\end{tabular}
\end{table}

% ══════════════════════════════════════════════════════════════════════════════
\section{Trade Implication}
% ══════════════════════════════════════════════════════════════════════════════

From the findings, there would exist a structural short with a short-medium time frame on Central American sovereign debt in this case. This could be done through sovereign CDS or relative value positions. With representative 10-year Central American sovereigns yielding approximately 6\% (duration $\sim$7 years), a 100 bps spread widening implies a $\sim$7\% mark-to-market price decline per \$100 face value --- opportunity for a potential trade value after fragmentation.

% ══════════════════════════════════════════════════════════════════════════════
\section{Limitations and Areas for Further Work}
% ══════════════════════════════════════════════════════════════════════════════

\textbf{Sample size.} 22 observations for a regression with 6 variables leaves only 15 degrees of freedom. A larger sample would improve identification, though data availability limits this in practice.

\textbf{Comparable quality.} Caribbean islands were the most direct comparable but still differ vastly from the hypothetical Central American nations along factors beyond simple geography. Within-region comparison is cleaner than global regression but still imperfect.

% ── Data Sources ───────────────────────────────────────────────────────────────
\vspace{8pt}
\hrule
\vspace{5pt}
{\small \textbf{Data Sources:} World Bank (GDP per capita, LPI); IMF WEO (debt-to-GDP, CA); INFORM Risk Index (disaster exposure); FAO Indicators (import dependency); TradingEconomics, FRED DGS10.}

\end{document}